\documentclass[11pt,a4paper]{article}
\usepackage[utf8]{inputenc}
\usepackage[T1]{fontenc}
\usepackage[italian]{babel}
\usepackage[margin=2cm]{geometry}
\usepackage{amsmath,amssymb}
\setlength{\parindent}{0pt}
\setlength{\parskip}{6pt}

\begin{document}

\begin{center}
\textbf{1SP --- Frazioni e numeri razionali}\\[4pt]
\large\textbf{Scheda 1 --- Equivalenza e proporzioni}
\end{center}

Componenti: \dotfill

\textbf{Tempo: 10 minuti.} Lavorate a coppie o in tre. Prima pensate individualmente per un minuto, poi confrontate i metodi. Scambiate chi scrive tra i due esercizi. Tutti devono saper spiegare i passaggi. Non usate la calcolatrice.

\section*{1. Trovale tutte}

Determinate \textbf{tutte} le frazioni equivalenti a
\[
-\frac{42}{63}
\]
che abbiano numeratore intero e denominatore intero positivo compreso tra 20 e 30, estremi inclusi.

Spiegate perché non ne mancano altre. Nella spiegazione collegate la proprietà invariantiva alla proporzionalità tra numeratore e denominatore.

\vspace{5.5cm}

\section*{2. Anche l'incognita può essere una frazione}

Riscrivete le seguenti proporzioni come uguaglianze tra frazioni e determinate l'incognita:
\[
\text{a)}\quad (-21):35=x:(-20),
\qquad
\text{b)}\quad 18:y=27:20.
\]

Non assumete che le incognite siano intere. Indicate le eventuali condizioni sui denominatori e verificate entrambe le risposte sostituendo i valori trovati. Per almeno una delle due, motivate l'uguaglianza mediante il denominatore comune o i prodotti incrociati.

\vfill
\textbf{Da discutere:} il rapporto resta costante quando si aggiunge lo stesso numero ai due termini di una frazione? Giustificate la risposta con un esempio.

\newpage

\begin{center}
\textbf{1SP --- Frazioni e numeri razionali}\\[4pt]
\large\textbf{Scheda 2 --- Espressioni e condizioni di validità}
\end{center}

Componenti: \dotfill

\textbf{Tempo: 10 minuti.} Lavorate a coppie o in tre, senza calcolatrice. Prima di calcolare controllate che le operazioni siano possibili. Scrivete i passaggi decisivi e scambiate chi scrive tra i due esercizi.

\section*{1. Un'espressione, più scelte da motivare}

Calcolate e riducete ai minimi termini:
\[
A=
\left[
\left(\frac56-\frac34\right):
\left(\frac79-\frac23\right)
\right]^{-1}
-\left(-\frac23\right)^{-2}.
\]

Indicate esplicitamente a quale numero si applica l'esponente $-1$. Spiegate il segno del termine $\left(-\frac23\right)^{-2}$ e controllate che ogni divisore sia non nullo.

\vspace{5cm}

\section*{2. Correggere non significa sempre trovare un risultato}

Uno studente svolge così:
\[
\begin{aligned}
B&=\frac{\frac23-\frac56}{\frac14-\frac13}
 :\left[\left(-\frac12\right)^{-2}-4\right]\\[6pt]
 &=2:[-4-4]\\[4pt]
 &=2:(-8)=-\frac14.
\end{aligned}
\]

Individuate il \textbf{primo passaggio sbagliato} e correggetelo. Stabilite quindi se l'espressione iniziale è definita: motivate la conclusione, anche se non potete assegnarle un valore.

\vfill
\textbf{Da discutere:} quali controlli avrebbero evitato l'errore prima di arrivare all'ultima riga?

\end{document}
